% ==============================================================================
% CAPÍTULO 5: CONCLUSÕES E TRABALHOS FUTUROS
% ==============================================================================

\section{CONCLUSÕES E TRABALHOS FUTUROS}
\label{sec:capitulo5_conclusoes}

Esta dissertação investigou, modelou e validou numericamente o acoplamento eletrotérmico, eletroquímico e radiativo orbital de um centro de dados espacial autônomo de $50\text{ kW}$ ($40\text{ kW}$ de carga útil de inteligência artificial em aceleradores FinFET de $5\text{ nm}$), operando no vácuo em órbita heliossíncrona terminadora (SSO $700\text{ km}$, LTAN 06h/18h), com análises comparativas em órbita terrestre baixa (LEO $550\text{ km}$, $53^\circ$) e geoestacionária (GEO $35.786\text{ km}$). A ausência de convecção atmosférica, a exposição ao espectro solar extraterrestre AM0 ($1361\text{ W/m}^2$), a degradação acumulativa por dose de deslocamento (DDD) e as transições abruptas de penumbra e umbra impõem restrições físicas severas que inviabilizam a simples transposição de arquiteturas terrestres de resfriamento e distribuição de potência para o espaço.

\subsection{Síntese dos Resultados e Fechamento dos Objetivos Específicos}
\label{sec:sintese_objetivos_cap5}

Os resultados analíticos e numéricos consolidados nos Capítulos~3 e~4 demonstram o cumprimento integral dos objetivos específicos estabelecidos na Seção~1.3.2:

\begin{enumerate}[label=\textbf{\arabic*.},leftmargin=*]
    \item \textbf{Arquitetura HVDC de $400\text{ V}$, Gerador Multijunção de 3 Diodos e Conversão GaN/DAB (Capítulo~3):} A adoção do barramento primário de $400\text{ V}$ em corrente contínua, combinada ao cabeamento em alumínio aeroespacial com \textit{derating} a vácuo, reduziu a massa de distribuição em $97{,}5\%$ ($522\text{ kg} \rightarrow 12{,}9\text{ kg}$) frente ao barramento legado de $28\text{ V}$, respeitando o limiar de descargas parciais por lei de Paschen com encapsulamento polimérico sólido e proteção SSPC em SiC com atuação ultrarrápida ($t_{trip} \le 15\ \mu\text{s}$, $\int i^2 dt = 145{,}5\text{ A}^2\text{s} \ll 1500\text{ A}^2\text{s}$). O gerador fotovoltaico de tripla junção $\text{GaInP}/\text{GaAs}/\text{Ge}$ ($158{,}19\text{ m}^2$, $52.416\text{ células}$ em $252\text{ strings} \times 208\text{ células}$ por satélite, totalizando $1.048.320\text{ células}$ e $3163{,}8\text{ m}^2$ na constelação de 20 nós), modelado com rigor pelo circuito equivalente de 3 diodos em série, garantiu $V_{mp} = 452{,}0\text{ V} > 400\text{ V}$ e razão cíclica $D = 0{,}885 \le 0{,}90$ no conversor Buck redutor no pior caso contínuo ($+65^\circ\text{C}$), assegurando potência de fim de vida (EOL, 10 anos em SSO $700\text{ km}$ sob degradação DDD e retenção de $90{,}42\%$) de $59{,}55\text{ kW}$ sobre os $65{,}55\text{ kW}$ BOL, superando com folga o requisito de $50{,}0\text{ kW}$ úteis por nó. A integração de diodos de bypass Schottky suprimiu o aquecimento reverso em \textit{hot spot} de $202{,}4^\circ\text{C}$ para seguros $60{,}0^\circ\text{C}$, enquanto o conversor Buck síncrono a $150\text{ kHz}$ com controle MPPT P\&O adaptativo ($f_{mppt} = 10\text{ kHz}$) manteve eficiência de rastreio de $99{,}92\%$ sob rampa térmica pós-eclipse de $15\text{ K/min}$. Nos conversores isolados DAB em GaN ($400\text{ V}/48\text{ V}$), o substrato de AlN com grafite pirolítico ($R_{th,total} = 1{,}12\text{ K/W}$) garantiu margem de $+136{,}6\%$ contra o limiar de fuga térmica ($2{,}65\text{ K/W}$), atingindo eficiência de $97{,}003\%$ a plena carga ($P_{loss,GaN} = 1235{,}7\text{ W}$ em 2 módulos ativos, ou $1025\text{ W}$ em compartilhamento $N+1$), com estabilidade de pequenos sinais assegurada pelo Critério de Middlebrook ($GM = 10$) sob cargas de potência constante (CPL).
    \item \textbf{Armazenamento Eletroquímico BESS 108S1P e Estimador Adaptativo AUKF Sage--Husa (Seção~3.5):} O banco de baterias $\text{LiNi}_{0{,}8}\text{Mn}_{0{,}1}\text{Co}_{0{,}1}\text{O}_2\text{ (NMC-811)}$ (NMC) de $9{,}0\text{ kWh}$ ($399{,}6\text{ V}$ nominais, pack 108S1P com 6 submódulos 18S), modelado por circuito equivalente de segunda ordem (\textit{Dual-Polarization}) acoplado à cinética de Arrhenius, demonstrou que a gestão ativa de carga pelo TEMS atende à norma ECSS-E-ST-20-32C, mantendo a profundidade de descarga contínua entre $20\%\text{ e }30\%$ em LEO para garantir retenção de capacidade superior a $80\%$ ao final de 10 anos, restringindo o consumo em sombra a $6{,}5\text{ kW}$ em LEO ($35{,}8\text{ min}$, $SoC_{min} = 50{,}46\%$), $18{,}0\text{ kW}$ em SSO ($18{,}8\text{ min}$, $SoC_{min} = 31{,}56\%$) e $4{,}5\text{ kW}$ em GEO ($72{,}0\text{ min}$, $SoC_{min} = 28{,}63\%$). O Filtro de Kalman Unscented Adaptativo (AUKF) com regulação de covariância de Sage--Husa livre de matrizes Jacobianas via fator de esquecimento $\lambda_{ad,k}$ sobre a dispersão dos pontos sigma convergiu em $< 90\text{ s}$ sob erro inicial severo de $30\%$ de $SoC$, mantendo erro médio quadrático de $0{,}037\%$ sob ruído severo de $150\text{ mV}_{rms}$, superando por $8\times$ a meta de projeto ($< 1{,}20\%$).
    \item \textbf{Acoplamento Radiativo Orbital, Rede Nodal de 5 Nós e Estabilidade de Lyapunov em FinFET de $5\text{ nm}$ (Capítulo~4):} A modelagem da dinâmica kepleriana com função de penumbra cônica analítica de classe $C^1$ e fatores de visão de calota esférica ECSS revelou que o radiador espacial bifacial em OSR sobre alma de grafite pirolítico encapsulado ($\text{APG}$, $k_{\parallel} = 1500\text{ W/(m}\cdot\text{K)}$, eficiência global de aleta $\eta_{fin} = 0{,}945$, área total $A_{rad} = 88{,}6\text{ m}^2$) rejeita os $41{,}984\text{ kW}$ térmicos internos em SSO $700\text{ km}$ mantendo $T_{rad,eq} = 52{,}8^\circ\text{C}$ e temperatura de junção $T_j = 77{,}8^\circ\text{C}$ em EOL ($\alpha_{s,EOL} = 0{,}14$, $\varepsilon_{ir,EOL} = 0{,}88$, $\Delta T_{j-rad} = 25{,}0\text{ K}$ com espalhamento bidimensional em câmara de vapor). A análise de sensibilidade de apontamento demonstrou início de \textit{throttling} em $|\delta\theta_{\text{thr}}| = 5{,}87^\circ$ e corte de emergência (\textit{thermal trip}, $T_j = 85^\circ\text{C}$) em $|\delta\theta_{\text{trip}}| = 15{,}24^\circ$. No domínio microeletrônico, a formulação de corrente de fuga sublimiar BSIM acoplada ao método direto de Lyapunov demonstrou que a operação em sobretensão ($V_{core} = 1{,}05\text{ V}$, $f_{clk} = 1{,}8\text{ GHz}$) sofre bifurcação sela-nó e fuga térmica irreversível em $T_{j,crit} = 83{,}4^\circ\text{C}$ ($\partial P_{leak}/\partial T_j > G_{th}$), justificando fisicamente o teto de projeto de $T_{j,max} = 80^\circ\text{C}$ ($V_{core} \le 0{,}85\text{ V}$).
    \item \textbf{Controle Preditivo TEMS-NMPC, Fadiga Termomecânica BGA e Eficácia Energética $PUE_{orbital}$ (Seções~4.5 e~4.6):} O controlador preditivo não linear baseado em modelo (TEMS-NMPC, horizonte $T_p = 60\text{ s}$) com antecipação de fila de tarefas (\textit{workload-aware pre-cooling}) suprimiu a oscilação de alta frequência (\textit{chattering}) dos controladores reativos ($450\text{ MHz/s} \rightarrow 65\text{ MHz/s}$), eliminou violações térmicas sob degraus de carga de $40\text{ kW}$ ($T_{j,max} = 79{,}4^\circ\text{C} < 80^\circ\text{C}$) e reduziu a amplitude de excursão térmica orbital $\Delta T_j$ na transição de eclipse em SSO de $22{,}5^\circ\text{C}$ para $7{,}0^\circ\text{C}$. Aplicando o modelo viscoplástico de Coffin-Manson-Engelmaier para soldas BGA $\text{SAC305}$ no pino de canto de dies de $28 \times 28\text{ mm}$ ($L_{DNP} = 19{,}80\text{ mm}$), comprovou-se que a operação sem TEMS em LEO entraria em colapso estrutural após $3.165\text{ ciclos}$ ($0{,}57\text{ ano}$, ou $6{,}8\text{ meses}$), ao passo que o TEMS-NMPC estende a vida útil por fadiga para $1{,}84 \times 10^5\text{ ciclos}$ ($34{,}53\text{ anos}$ em SSO, $+245{,}3\%$ sobre a meta de 10 anos). Globalmente, o centro de dados orbital alcançou $PUE_{orbital} = 1{,}050$ na fronteira térmica ($1{,}057$ na fronteira elétrica global com 2 módulos DAB ativos) e rendimento computacional específico de $\eta_{comp} = 0{,}476\text{ TOPS/W}$ ($20.000\text{ TOPS}$ em 80 dies a $1{,}8\text{ GHz}$), superando com folga instalações terrestres de alto desempenho.
\end{enumerate}

\subsection{Limitações Epistêmicas e Condições para Fechamento Experimental}
\label{sec:limitacoes_epistemicas_cap5}

Em consonância com o rigor metodológico de engenharia aeroespacial adotado neste trabalho, cumpre explicitar as fronteiras de validade da verificação \textit{in silico} realizada:
\begin{itemize}[leftmargin=*]
    \item \textbf{Calibração Empírica de Parâmetros concentrados e Termomecânicos:} Os coeficientes de degradação óptica ($\alpha_{s,EOL}$, $\varepsilon_{ir,EOL}$), as resistências equivalentes da rede térmica de 5 nós ($R_{12}$ a $R_{45}$), a constante de ductilidade efetiva em fadiga ($\text{SAC305}$ sob encapsulamento rígido e \textit{underfill}) e os parâmetros cinéticos do BESS (NMC 108S1P de $9{,}0\text{ kWh}$) foram ancorados na literatura aberta e em folhas de dados de qualificação aeroespacial, requerendo campanha de ensaios de ciclagem térmica em câmara de termovácuo (TVAC) e caracterização calorimétrica de lote de voo para fechamento definitivo de missão (nível TRL 6--8).
    \item \textbf{Acoplamento Termo-Hidráulico Bifásico Unidimensional:} O circuito bombeado bifásico (MPFL/LHP) foi representado por sua condutância global efetiva e limiar paramétrico de secagem capilar (\textit{dry-out}). Fenômenos transitórios de instabilidade de ebulição nucleada, ondas de choque de pressão vapor-líquido sob microgravidade e distribuição de fluxo em múltiplos evaporadores paralelos demandam modelagem fluidodinâmica bifásica distribuída (1D/3D) e testes em bancada suborbital ou na Estação Espacial Internacional.
\end{itemize}

\subsection{Recomendações para Trabalhos Futuros}
\label{sec:trabalhos_futuros_cap5}

Como desdobramentos diretos desta pesquisa para a consolidação de infraestruturas orbitais de processamento de inteligência artificial, propõem-se quatro linhas de investigação:
\begin{enumerate}[label=\textbf{(\alph*)},leftmargin=*]
    \item \textbf{Implementação \textit{Hardware-in-the-Loop} (HIL) em FPGA Tolerante à Radiação:} Síntese em ponto fixo e validação em tempo real do estimador AUKF Sage--Husa e do solver quadrático sequencial do TEMS-NMPC em arquiteturas FPGA espaciais (ex.: Xilinx Virtex-5QV / Kintex UltraScale XQRKU060), acopladas a um emulador eletrônico do barramento HVDC de $400\text{ V}$ e de células de bateria reais em câmara climática.
    \item \textbf{Escalonamento Federado de Cargas de IA em Constelações via OISL:} Extensão do controlador TEMS-NMPC local de um único satélite para um protocolo de controle preditivo distribuído em constelação (\textit{Distributed MPC}), roteando dinamicamente lotes de inferência e gradientes de treinamento através dos enlaces ópticos inter-satélites ($400\text{ Gbps}$) dos nós que ingressam em penumbra ou sofrem desvios de atitude para nós iluminados com folga térmica nos radiadores.
    \item \textbf{Tolerância a Falhas por Radiação Ionizante (\textit{Single-Event Effects}) no Co-Design Eletrotérmico:} Integração estocástica de taxas de \textit{Single-Event Upset} (SEU) e \textit{Single-Event Latchup} (SEL) induzidas pela Anomalia Magnética do Atlântico Sul (SAA) e por prótons solares diretamente na função custo do TEMS, combinando regulação de tensão/frequência (DVFS) com sobrecarga energética de verificação de paridade, \textit{checkpointing} em memória HBM3 com código corretor de erros (ECC) e redundância modular seletiva.
    \item \textbf{Bancada Experimental de Microcanal Bifásico e Substratos Compósitos AlN/APG em Vácuo:} Construção e ensaio experimental em câmara de alto vácuo ($< 10^{-5}\text{ Pa}$) de um protótipo em escala reduzida ($1\text{ a }2\text{ kW}$) integrando aquecedores planares de alta densidade superficial ($> 60\text{ W/cm}^2$), conversores GaN sobre substrato de nitreto de alumínio e dissipadores pirolíticos, visando levantar experimentalmente as curvas de bifurcação de Lyapunov e validar os limiares da Matriz FMECA consolidada.
\end{enumerate}
